% page 435 du fac-simile (image p-434 du scan)
% bloc 165.1 x 242.0 mm ; marge gauche 36.9 mm
\pgc{15.240mm}{0.000mm}{
\l{\cel{56}427}
\l{}
\l{\sou{patogenezo}. (\sou{patol.)}\cel{2}Exploro di la mekanismo per qua la kauzi}
\l{\cel{3}morbifera efikas a la organismo vivoza por maladigar olca.}
\l{}
\l{\sou{patologio}. Parto di medicino, qua studias la morbi, traktata}
\l{\cel{3}kom tala. - DEFIRS.}
\l{}
\l{\sou{patoso}. (\sou{retoriko})\cel{2}Uzo exajerita di la moyeni apta agitar forte}
\l{\cel{3}la pasioni. - DEFR.}
\l{}
\l{\sou{patro}. I. La persono qua genitis filio. - II. (\cel{1}\sou{fig.}) Titulo ye}
\l{\cel{3}veneraceso, respekteso. - III. Membro di kongregaciono reli-}
\l{\cel{3}giala (katolika). - EFIS.}
\l{}
\l{\sou{patrio}. Lando ube onu naskis ed a qua onu apartenas kom civitano.}
\l{\cel{3}- FIS.}
\l{}
\l{\sou{patriarko}.- I. (\sou{biblo}). Familio-chefo di qua la vivo esis tre}
\l{\cel{3}longa. - II. (\sou{metaf}.)Oldo veneracinda, cirkondata da multa}
\l{\cel{3}decendanti. - III. (\sou{historio di la eklezio kristana}).Episkopo}
\l{\cel{3}di la episkopeyo komencala. - IV. Chefo di la eklezio Greka.}
\l{\cel{3}- DEFIRS.}
\l{}
\l{\sou{patrico}. (en\cel{1}\sou{Roma, antique}). Oficiero alta-klasa, en la en la}
\l{\cel{3}periodi +ultima di la imperio Romana. - DEFIRS.}
\l{}
\l{\sou{patrimonio}.\cel{2}Havajo, domeno, quan onu heredis de la ancestri}
\l{\cel{3}patrulala o patrinala. - DEFIRS.}
\l{}
\l{\sou{patriota}.\cel{2}Devota por sua patrio. - DEFIRS.}
\l{}
\l{\sou{patristiko.}\cel{2}(\sou{religio katolika}).\cel{2}Cienco pri la doktrino di la}
\l{\cel{3}patri ekleziala. -}
\l{}
\l{\sou{patroliar}. (\sou{netrans.}) (\sou{aludante detachmento)}. Cirkular por la}
\l{\cel{3}sekureso di urbo, di kampeyo, por verifikar kad omno esas en}
\l{\cel{3}bona stando. - DEFIRS.}
\l{}
\l{\sou{patrologio}. (\sou{eklezio katolika}).Cienco qua traktas la verki, la}
\l{\cel{3}biografio di la patri ekleziala. - DEFIS.}
\l{}
\l{\sou{patrono.}\cel{1}I. (en\cel{1}\sou{Roma, dum la epo}ki\cel{1}\sou{antiqua}). Persono richa,}
\l{\cel{3}potenta,\cel{2}cirkum qua su grupigis klienti. - II. La mastro}
\l{\cel{3}di establisuro komercala, industriala, che qua salario-laboras}
\l{\cel{3}kontor-employati, fabrik-employati. - DEFIRS.}
\l{}
\l{pauperismo. La kalamitato \textquotedbl{}povreso\textquotedbl{}, en stato. - DEFRS.}
\l{}
\l{pauzar. (netrans.) I. Interruptar la agado. - II. (muziko).}
\l{\cel{3}Cesar la kanto,la pleo instrumentala, dure di noto kompleta.}
\l{\cel{3}- DEFIRS.}
\l{}
\l{\cel{1}\sou{pavo}. Singla de la bloki de greso, granito, e c., kubigita per}
\l{\cel{4}talio, quin onu asemblas por igar plu rezistiva la sulo di}
\l{\cel{3}choseo,di voyo,di strado, di voyeto. - EF.}
}
